%% Chapter 5: Conclusion & Future Work (chapters/chapter05.tex)

% ==============================================================================
% CONTENT BLUEPRINT — Chapter 5: Conclusion and Future Work
% ==============================================================================
% Reading audience: anyone skimming the project to decide if the work is solid.
% Goal: in 5–8 pages, deliver the verdict on each objective and the
% roadmap ahead.
%
% ----------------------------------------------------------------------------
% §5.1 Summary of Work and Findings  (2–3 pages)
%   • Start with a single paragraph that, in plain language, states:
%       – What was studied (one sentence),
%       – How it was studied (one sentence),
%       – What was found (one sentence with the headline number).
%   • Then a bullet list of the headline findings, one per objective in
%     §1.4. Each bullet states the finding, the supporting figure (§3.x) and
%     citation (if any).
%   • Restate explicitly which objectives were met, which were partially
%     met, and which were not — and why.
%
% §5.2 Contributions of the Project  (½–1 page)
%   • Enumerate the deliverables:
%       – Code (with public repo URL and DOI),
%       – Dataset (Zenodo DOI),
%       – Calibration constants / values,
%       – Methodological workflow / SOPs.
%   • This is what can be reused by the next student.
%
% §5.3 Limitations  (½–1 page)
%   • Restate the principal limitations from §3.4 and §4.3.
%   • Be concise. The committee will not accept a six-page criticism of your
%     own work; one paragraph with a bulleted breakdown is ideal.
%
% §5.4 Recommended Future Work  (1–2 pages)
%   • 3–5 directions, each one paragraph:
%       – What would be done,
%       – What new physics / engineering problem it unlocks,
%       – What would be required in terms of equipment / time / funds.
%   • Each item should be specific enough that another student could pick
%     one up and turn it into their own project proposal.
%
% §5.5 Closing Statement  (¼ page, optional)
%   • One paragraph. Reflect on the broader significance of the work or the
%     author's learning. Keep it professional — avoid personal reflections
%     about feelings or family.
%
% Writing-style rules for this chapter:
%   • Use declarative, past tense. The work is finished.
%   • Each numerical claim must reference either §3 or §4 (or be a repeat
%     of §1's headline objective).
%   • No new figures. If you find you need a new figure, return to Chapter 3.
% ==============================================================================

\chapter{Conclusion and Future Work}
\label{ch:conclusion}

This B.Sc. project template has demonstrated the design, configuration, and formatting requirements of a standard Physics report at Shahjalal University of Science and Technology.

\section{Summary of Contributions}
We have successfully implemented:
\begin{enumerate}
    \item A structural template for writing B.Sc. Physics projects.
    \item Formatting guidelines for equations, figures, tables, and listings.
    \item Sample TikZ block diagrams and optical table layout scripts.
    \item Validation plots using pgfplots to minimize dependency on external tools.
\end{enumerate}

\section{Proposed Future Investigations}
For future revisions, we propose:
\begin{itemize}
    \item Developing a Python GUI wrapper to read data directly from digital power meters.
    \item Adding support for double-pass Z-scan geometries.
    \item Automating build procedures via GitHub actions or containerized TeX compilers.
\end{itemize}
This boilerplate forms a robust launchpad for academic writing within the physics department.
